\begin{tikzpicture}[
  kolom/.style={font=\scriptsize\bfseries, align=center, text width=1.28cm},
  baris/.style={font=\scriptsize, align=left, text width=3.1cm, anchor=west},
  pemutus/.style={circle, fill=praditagreen, draw=black, minimum size=3mm,
    inner sep=0pt},
  masukan/.style={circle, fill=white, draw=black, thick, minimum size=3mm,
    inner sep=0pt},
  kini/.style={rectangle, fill=red!22, draw=black, minimum size=2.8mm,
    inner sep=0pt}
]
  % Kisi 6 kolom arketipe dan 5 baris keputusan; lebar total 11,7 cm.
  \fill[praditagreen!15] (0, 0) rectangle (11.7, 0.85);
  \draw[gray] (0, -3.5) grid[xstep=11.7, ystep=0.7] (11.7, 0);
  \foreach \posisi in {3.3, 4.7, 6.1, 7.5, 8.9, 10.3, 11.7} {
    \draw[gray] (\posisi, -3.5) -- (\posisi, 0.85);
  }
  \draw[gray] (0, -3.5) rectangle (11.7, 0.85);
  \node[baris, font=\scriptsize\bfseries] at (0, 0.425) {Keputusan};
  \node[kolom] at (4.0, 0.425) {Monarki bisnis};
  \node[kolom] at (5.4, 0.425) {Monarki TI};
  \node[kolom] at (6.8, 0.425) {Feodal};
  \node[kolom] at (8.2, 0.425) {Federal};
  \node[kolom] at (9.6, 0.425) {Duopoli TI};
  \node[kolom] at (11.0, 0.425) {Anarki};
  \node[baris] at (0, -0.35) {Prinsip TI};
  \node[baris] at (0, -1.05) {Arsitektur TI};
  \node[baris] at (0, -1.75) {Infrastruktur TI};
  \node[baris] at (0, -2.45) {Kebutuhan aplikasi bisnis};
  \node[baris] at (0, -3.15) {Investasi dan prioritas TI};
  % Prinsip TI
  \node[pemutus] at (4.0, -0.35) {};
  \node[masukan] at (8.2, -0.35) {};
  \node[kini] at (5.4, -0.35) {};
  % Arsitektur TI
  \node[pemutus] at (5.4, -1.05) {};
  \node[masukan] at (8.2, -1.05) {};
  \node[kini] at (6.8, -1.05) {};
  % Infrastruktur TI
  \node[pemutus] at (5.4, -1.75) {};
  \node[masukan] at (9.6, -1.75) {};
  \node[kini] at (6.8, -1.75) {};
  % Kebutuhan aplikasi bisnis
  \node[pemutus] at (8.2, -2.45) {};
  \node[masukan] at (6.55, -2.45) {};
  \node[kini] at (7.05, -2.45) {};
  % Investasi dan prioritas TI
  \node[pemutus] at (4.0, -3.15) {};
  \node[masukan] at (9.6, -3.15) {};
  \node[kini] at (5.4, -3.15) {};
  % Keterangan penanda
  \node[pemutus] at (0.2, -3.95) {};
  \node[font=\scriptsize, anchor=west] at (0.4, -3.95) {pemutus usulan};
  \node[masukan] at (3.6, -3.95) {};
  \node[font=\scriptsize, anchor=west] at (3.8, -3.95) {pemberi masukan usulan};
  \node[kini] at (7.9, -3.95) {};
  \node[font=\scriptsize, anchor=west] at (8.1, -3.95) {pemutus saat ini};
\end{tikzpicture}
